\documentclass[11pt]{article}
\usepackage[utf8]{inputenc}
\usepackage[T1]{fontenc}
\usepackage{geometry}
\usepackage{amsmath, amssymb}
\geometry{a4paper, margin=2cm}

\begin{document}

\begin{center}
{\LARGE \textbf{Gabarito - Aprofundamento em Matemática: \\  Análise Combinatória}}
\end{center}

\bigskip

\textbf{1.} Uma lanchonete oferece 3 tipos de sanduíche e 2 tipos de suco. Quantas combinações diferentes podem ser escolhidas?

\textbf{Resposta:}\\
Multiplicação de opções: $3 \times 2 = 6$. \\
\textbf{Resposta: 6.}

\bigskip

\textbf{2.} Um código possui 2 letras e 3 números, com repetição permitida. Quantos códigos diferentes podem ser formados?

\textbf{Resposta:}\\
Letras: $26$ possibilidades cada.  
Números: $10$ possibilidades cada.  
Total: $26 \times 26 \times 10 \times 10 \times 10 = 26^2 \cdot 1000$. \\
\textbf{Resposta: } $26 \times 26 \times 1000$.

\bigskip

\textbf{3.} Quantos números de 3 algarismos distintos podem ser formados com \{1,2,3,4\}?

\textbf{Resposta:}\\
Primeiro algarismo: 4 opções.  
Segundo: 3 opções.  
Terceiro: 2 opções.  
Total: $4 \times 3 \times 2 = 24$. \\
\textbf{Resposta: 24.}

\bigskip

\textbf{4.} De quantas maneiras diferentes podemos organizar 3 livros distintos?

\textbf{Resposta:}\\
Permutação de 3: $3! = 6$. \\
\textbf{Resposta: 6.}

\bigskip

\textbf{5.} De quantas maneiras podemos escolher 2 alunos entre 5?

\textbf{Resposta:}\\
Combinação: $\binom{5}{2} = \frac{5 \cdot 4}{2} = 10$. \\
\textbf{Resposta: 10.}

\bigskip

\textbf{6.} Quantas equipes de 3 jogadores podem ser formadas a partir de 6 pessoas?

\textbf{Resposta:}\\
$\binom{6}{3} = \frac{6 \cdot 5 \cdot 4}{6} = 20$. \\
\textbf{Resposta: 20.}

\bigskip

\textbf{7.} Uma pessoa possui 4 camisas e 3 calças. Quantos trajes diferentes podem ser escolhidos?

\textbf{Resposta:}\\
$4 \times 3 = 12$. \\
\textbf{Resposta: 12.}

\bigskip

\textbf{8.} Quantos anagramas podem ser feitos com “SOL”?

\textbf{Resposta:}\\
Permutação de 3 letras distintas: $3! = 6$. \\
\textbf{Resposta: 6.}

\bigskip

\textbf{9.} Um comitê de 2 pessoas será escolhido entre 4 candidatos. Quantas maneiras existem?

\textbf{Resposta:}\\
$\binom{4}{2} = \frac{4 \cdot 3}{2} = 6$. \\
\textbf{Resposta: 6.}

\bigskip

\textbf{10.} Um menu oferece 2 entradas, 3 pratos principais e 2 sobremesas. Quantas refeições completas diferentes existem?

\textbf{Resposta:}\\
Multiplicação das escolhas:  
\[
2 \times 3 \times 2 = 12.
\]
\textbf{Resposta: 12.}

\end{document}

